% Material removed from main.tex in the round-6 simplification (2026-11-27).
% Kept verbatim from commit 5d2a07b for reference; it is not part of the paper and is not compiled.
% Line numbers refer to main.tex at 5d2a07b.

% ---- Abstract and intro claims (lines 17-22, 45-49, 58-61) ----
$d^2\cdot d!/\min_\Omega h_f$. Harris's and Janson's inequalities, applied to cells whose size tends to $0$ for fixed $n$, give exact
Poisson bounds, and for $j\le d-2$ a restriction of the spanned flats (heuristically needed when $j<d/2$) makes the dependence through shared points negligible. The
deepest uncovered points converge to the minimizers of $h_f$. We compute the limit for balls, e.g.\ $54$ for planes and $10$ for
lines in the unit $3$-ball, and for the square, and show that among rotation-invariant densities on a ball the uniform one is
optimal only in dimension $3$. As $n$ points can be placed so that every point of $\Omega$ is within $O(n^{-d})$ of a spanned
hyperplane (a companion paper), random points lose exactly a factor of order $\log n$.

be any compact set that is the closure of its interior. For hyperplanes and uniform points the limit is $c_\Omega=d^2\cdot
d!/\min_\Omega h_\Omega$, so with Bang's bound $\rho^{(n)}/\rho^*_n\le(1+o(1))\frac{2c_\Omega}{d!\,w(\Omega)}\log n$ almost surely (for
instance $\frac{9\pi^2}{16}\log n\approx5.55\log n$ for the disk and $\approx6.74\log n$ for the square), while $\rho^*_n=O(n^{-d})$ gives
$\rho^{(n)}/\rho^*_n\ge c\log n$: random points lose exactly a factor of order $\log n$ (and $(\log n)^{1/(d-j)}$ for $j$-flats in
the ball).

(Corollary~\ref{cor:where}), e.g.\ to the edge midpoints of a square. (iv) By the affine Blaschke--Petkantschin formula
\cite[Theorem~7.2.7]{SW}, $h_f(z)$ is an integral over the sections of $\Omega$ by flats through $z$ (Figure~\ref{fig:lines}). This gives the limit for
balls and every $j$, e.g.\ $54$ for planes and $10$ for lines in the unit $3$-ball (Corollary~\ref{cor:ball}); for hyperplanes
in the plane, a cubic moment of the chord lengths through $z$, the minimizers for centrally symmetric bodies and the square, and
the fact that among rotation-invariant densities on a ball the uniform one is optimal in dimension $3$ but not otherwise
(Section~\ref{sec:plane}).

% ---- Corollary 2.3 for beta-densities and further ball constants (lines 141-153) ----
\begin{corollary}\label{cor:ball}
Let $\Omega=\ball^d$ be the unit ball, $\beta\ge0$, $f_\beta$ the density on $\ball^d$ proportional to $(1-|x|^2)^\beta$ ($f_0$ is uniform),
$k=\frac{j(d+1)}2+\beta(j+1)$, and $\mathrm B$ the Beta function. Then $h_{f_\beta}(z)$ depends only on $|z|$ and is strictly decreasing in
$|z|\in[0,1]$, so its minimizers on $\ball^d$ are the points with $|z|=1$, and
\[
  c_{f_\beta}=\frac{(j+1)!\,d(j+1)}{2}\cdot\frac{\mathrm B(\frac a2,k+1)\,\mathrm B(\frac a2,\frac j2)}{\mathrm B(\frac a2,k+\frac j2)}.
\]
For hyperplanes, with $J(q)=\int_{-1}^1(1-u^2)^q\,du=\mathrm B(\frac12,q+1)=\sqrt\pi\,\Gamma(q+1)/\Gamma(q+\frac32)$, this is
$c_{f_\beta}=\frac{d^2\cdot d!}{2}J(\frac{d-3}2)J(k)/J(k+\frac{d-3}2)$ with $k=\frac{d^2-1}2+d\beta$; for uniform points $c_{\ball^2}=9\pi^2/8\approx11.10$,
$c_{\ball^3}=54$, $c_{\ball^4}=2111870475\pi^2/67108864\approx310.59$ and $c_{\ball^5}=27000/13$. For uniform points and lines ($j=1$) the
constant is $10$ in $\ball^3$, $525\pi^2/512\approx10.12$ in $\ball^4$ and $21/2$ in $\ball^5$; for planes in $\ball^4$ and $\ball^5$ it is $36$
and $30$.
\end{corollary}

% ---- Proof of Corollary 2.3, beta case (lines 547-566) ----
\begin{proof}[Proof of Corollary~\ref{cor:ball}]
Let $f=f_\beta=\alpha_\beta(1-|x|^2)^\beta$ on $\ball^d$. For $L\in G(d,j)$ and $t\in L^\perp$ with $|t|<1$, $(L+t)\cap\ball^d$ is the $j$-ball of
radius $\varrho=(1-|t|^2)^{1/2}$ about $t$, on which $f(t+\ell)=\alpha_\beta(\varrho^2-|\ell|^2)^\beta$; substituting $\ell_i=\varrho w_i$ in the
definition of $G_L(t)$ gives $G_L(t)=\alpha_\beta^{j+1}\Psi_\beta\,\varrho^{j(j+1)+ja+2\beta(j+1)}=\alpha_\beta^{j+1}\Psi_\beta(1-|t|^2)_+^{k}$, where
$\Psi_\beta=\int_{(B^j)^{j+1}}\prod_i(1-|w_i|^2)^\beta\vol_j(\operatorname{conv}\{w_i\})^a\,dw$ over the unit ball $B^j$ of $L$ (independent of $L$ by
isometry invariance). By \eqref{eq:BPj} with all flats, $1=b\alpha_\beta^{j+1}\Psi_\beta\int_{\R^a}(1-|t|^2)_+^k\,dt=b\alpha_\beta^{j+1}\Psi_\beta\frac{|\Sph^{a-1}|}2\mathrm B(\frac a2,k+1)$.
(Blaschke--Petkantschin computations of this kind for beta-distributed points in the ball are carried out in \cite{KTT};
see also \cite{RM}.) By Lemma~\ref{lem:ig}(c) and $\kappa_a=|\Sph^{a-1}|/a$, for $z\in\R^d$,
\[
  h_f(z)=b\kappa_a\alpha_\beta^{j+1}\Psi_\beta\,\E\bigl(1-|z|^2\mathcal Y\bigr)_+^k=\frac{2}{a\,\mathrm B(\frac a2,k+1)}\E\bigl(1-|z|^2\mathcal Y\bigr)_+^k,
\]
where $\mathcal Y=|\pi_Le|^2$ for $\nu_j$-uniform $L$ and a fixed unit vector $e$, which has the Beta$(\frac a2,\frac j2)$ law (proof of
Lemma~\ref{lem:ig}(b)); so $h_f(z)$ depends only on $|z|$. Since $k>0$, $0<\mathcal Y<1$ almost surely and so $1-|z|^2\mathcal Y>0$ for $|z|\le1$, $h_f$ is strictly decreasing in
$|z|\in[0,1]$, and as $\E(1-\mathcal Y)^k=\mathrm B(\frac a2,\frac j2+k)/\mathrm B(\frac a2,\frac j2)$,
$m_f=2\mathrm B(\frac a2,k+\frac j2)/(a\,\mathrm B(\frac a2,k+1)\mathrm B(\frac a2,\frac j2))$; so $c_f=(j+1)!\,d(j+1)/(am_f)$ (Theorem~\ref{thm:main}) is the
stated formula. For $a=1$, $\mathrm B(\frac12,q+1)=J(q)$ (substitute $x=u^2$), which gives the hyperplane form. The values follow
from $J(0)=2$, $J(-\frac12)=\pi$, $J(q)=\frac{2q}{2q+1}J(q-1)$ and $\mathrm B(x,y)=\Gamma(x)\Gamma(y)/\Gamma(x+y)$; e.g.\ for $d=3$ and
hyperplanes, $J(4)$ cancels and $c_{\ball^3}=27J(0)=54$; for lines in $\ball^3$, $a=2$, $k=2$ and
$c=6\,\mathrm B(1,3)\mathrm B(1,\frac12)/\mathrm B(1,\frac52)=6\cdot\frac13\cdot2\cdot\frac52=10$. (Exact checks: \texttt{ball\_flats.py} in \cite{Repo}.)
\end{proof}

% ---- Hyperplane normal coordinates, Psi, Z_Omega, eq. (BP) (lines 214-252, 267-268, 292-294) ----
$\{(L,t):t\in L^\perp,\ L+t\in\mathcal A\}$ is Borel. For hyperplanes we also use the normal coordinates: for $\nu\in\Sph^{d-1}$ (the unit
sphere, with surface measure $d\nu$) and $p\in\R$ let $H(\nu,p)=\{x:\ip x\nu=p\}$ and $G_\nu(p)=G_{\nu^\perp}(p\nu)$, the integral of
$\prod_if(x_i)\vol_{d-1}(\operatorname{conv}\{x_i\})$ over $d$-tuples in $H(\nu,p)$. For a compact convex set $K$ in a hyperplane let
\[
  \Psi(K)=\int_{K^d}\vol_{d-1}(\operatorname{conv}\{x_1,\dots,x_d\})\,d\sigma(x_1)\cdots d\sigma(x_d),
\]
and $\Psi(\emptyset)=0$; $\Psi$ is invariant under isometries, monotone and homogeneous of degree $d^2-1$; let $\Psi_{d-1}>0$ be $\Psi$
of a unit $(d-1)$-ball. For the uniform density, $G_\nu(p)=\Psi(\Omega\cap H(\nu,p))/|\Omega|^d$; let
$Z_\Omega=\int_{\Sph^{d-1}}\int_\R\Psi(\Omega\cap H(\nu,p))\,dp\,d\nu$.

\begin{lemma}[Representation]\label{lem:rep}
$(L,t)\mapsto G_L(t)$ is Borel, $G_L\le G_{\max}:=f_+^{j+1}(\kappa_j\Delta^j)^{j+1+a}$, and there is $b>0$ depending only on $d$ and
$j$ such that for measurable $\mathcal A$,
\begin{equation}\label{eq:BPj}
  \Pr(F\in\mathcal A)=b\int_{G(d,j)}\int_{L^\perp}\Cone\{L+t\in\mathcal A\}\,G_L(t)\,dt\,\nu_j(dL).
\end{equation}
For hyperplanes, with $b_d=b/|\Sph^{d-1}|$,
\begin{equation}\label{eq:BP}
  \Pr(H\in\mathcal A)=b_d\int_{\Sph^{d-1}}\int_\R\Cone\{H(\nu,p)\in\mathcal A\}\,G_\nu(p)\,dp\,d\nu ,
\end{equation}
and $|\Omega|^d=b_dZ_\Omega$.
\end{lemma}

\begin{proof}
$G$ is Borel by Appendix~\ref{app:meas}. The
integrand of $G_L(t)$ vanishes unless every $\ell_i$ lies in $K=\{\ell\in L:\ell+t\in\Omega\}$, a convex set of diameter $\le\Delta$, so
$\vol_jK\le\kappa_j\Delta^j$ and $\vol_j(\operatorname{conv}\{\ell_i\})\le\vol_jK$; this gives the bound. The affine Blaschke--Petkantschin
formula \cite[Theorem~7.2.7]{SW} with $q=j$ states, for measurable $\varphi\ge0$ on $(\R^d)^{j+1}$,
\[
  \int_{(\R^d)^{j+1}}\varphi\,dx=b_{dj}(j!)^a\int_{G(d,j)}\int_{L^\perp}\int_{(L+t)^{j+1}}\varphi\,\vol_j(\operatorname{conv}\{x_i\})^a\,dx\,dt\,\nu_j(dL),
\]
with $b_{dj}>0$, the measure on $j$-flats being normalized in \cite{SW} as the integral over $L\in G(d,j)$ (with $\nu_j$) of
Lebesgue measure on $L^\perp$. Take $\varphi(x)=\prod_if(x_i)\Cone\{x\in\mathcal G_j,\ F(x)\in\mathcal A\}$ (Borel by Appendix~\ref{app:meas}); its integral is $\Pr(F\in\mathcal A)$ as the complement of $\mathcal G_j$
is null. On $(L+t)^{j+1}$ every tuple of positive weight spans $L+t$, and $x_i=\ell_i+t$ gives \eqref{eq:BPj} with $b=b_{dj}(j!)^a$.
For $j=d-1$, $\nu_{d-1}$ is the image of $d\nu/|\Sph^{d-1}|$ under $\nu\mapsto\nu^\perp$ (by uniqueness of the rotation-invariant
probability measure), and $\int_{\nu^\perp{}^\perp}g(t)\,dt=\int_\R g(p\nu)\,dp$; this gives \eqref{eq:BP} (each hyperplane appears as
$(\nu,p)$ and $(-\nu,-p)$, but the inner integral depends only on $\nu^\perp$, so no factor $\frac12$ arises). The uniform density
with $\mathcal A$ all hyperplanes gives $|\Omega|^d=b_dZ_\Omega$. (The value of $b$ is never taken from \cite{SW}; where it is needed
it is determined by $|\Omega|^d=b_dZ_\Omega$.)
(c) As $t\to0$, $t^{-a}\Pr(\dist(z,F)<t)\to h_f(z)=b\kappa_a\int_{G(d,j)}G_L(\pi_Lz)\,\nu_j(dL)$ uniformly in $z\in\R^d$; $h_f$ is
continuous and positive on $\R^d$, and $h_f\ge c_1>0$ on $\Omega$. For hyperplanes, $h_f(z)=2b_d\int_{\Sph^{d-1}}G_\nu(\ip z\nu)\,d\nu$, and
for the uniform density $h_\Omega(z)=\frac2{Z_\Omega}\int_{\Sph^{d-1}}\Psi(\Omega\cap H(\nu,\ip z\nu))\,d\nu$.
$h_f(z)\ge b\kappa_af_-^{j+1}(\sqrt3r_0/2)^{j(d+1)}\Psi_{j,a}c's_0^a>0$, and for $z\in\Omega$ ($R=\Delta$) this is a constant $c_1$. For
hyperplanes, $\pi_{\nu^\perp}z=\ip z\nu\nu$ and $\kappa_1=2$, and the image of $d\nu/|\Sph^{d-1}|$ is $\nu_{d-1}$; with $b_d=b/|\Sph^{d-1}|$ and
$G_\nu=\Psi(\Omega\cap H(\nu,\cdot))/|\Omega|^d$ for the uniform density, this gives the last two formulas.

% ---- Proposition 7.2 (radial densities) with proof (lines 592-633) ----
\begin{proposition}\label{prop:radial}
Let $\Omega=\ball^d$, let $f$ be rotation invariant, and consider hyperplanes.
(a) If $d=3$, then $h_f(z)=\Pr(\dist(0,H)\le|z|)/|z|$ for $0<|z|\le1$. Hence $c_f\ge54=c_{\ball^3}$, with equality if and only
if $\Pr(\dist(0,H)\le r)\ge r$ for $0<r\le1$; in particular if $\dist(0,H)$ has a nonincreasing density on $[0,1]$, as for every
$f_\beta$, $\beta\ge0$.
(b) As a function of $\beta\ge0$, $c_{f_\beta}$ is strictly increasing if $d=2$, constant if $d=3$, and strictly decreasing if
$d\ge4$.
(c) If $d=2$ and $f(x)=(1+2|x|^2)/(2\pi)$, then $h_f$ is nonincreasing in $|z|$ and $c_f=1225\pi^2/1138\approx10.62<c_{\ball^2}$.

\noindent Thus among continuous rotation-invariant densities on the ball, the uniform density minimizes $c_f$ if $d=3$
(though not uniquely: every $f_\beta$ attains $54$), but not if
$d=2$ (mass pushed outwards helps) or $d\ge4$ (mass pulled inwards helps).
\end{proposition}

\begin{proof}[Proof of Proposition~\ref{prop:radial}]
(a) By rotation invariance $G_\nu(p)=g(p)$ does not depend on $\nu$, and $g$ is even and vanishes for $|p|>1$. By \eqref{eq:BP},
$\Phi(r)=\Pr(\dist(0,H)\le r)=b_3\int_{\Sph^2}\int_{-r}^rg(p)\,dp\,d\nu=4\pi b_3\int_{-r}^rg(p)\,dp$, so $\dist(0,H)$ has the density
$8\pi b_3g$ on $[0,1]$, and $\Phi(1)=1$. For $z\ne0$ and $\nu$ uniform on $\Sph^2$, $\ip z\nu$ is uniform on $[-|z|,|z|]$ (Archimedes);
so by Lemma~\ref{lem:ig}(c), $h_f(z)=2b_3\cdot4\pi\,\E g(\ip z\nu)=\frac{4\pi b_3}{|z|}\int_{-|z|}^{|z|}g(p)\,dp=\Phi(|z|)/|z|$. By continuity
of $h_f$, $m_f=\inf_{0<r\le1}\Phi(r)/r\le \Phi(1)=1$, with equality if and only if $\Phi(r)\ge r$ on $(0,1]$; so $c_f=d^2\cdot d!/m_f\ge54$,
with the same condition for equality. If $\Phi$ has a nonincreasing density, $\Phi$ is concave with $\Phi(0)=0$, so $\Phi(r)\ge r\Phi(1)=r$. For $f_\beta$,
$g$ is proportional to $(1-p^2)_+^{k}$, $k=\frac{d^2-1}2+d\beta$, by the proof of Corollary~\ref{cor:ball}, which is nonincreasing on $[0,1]$.

(b) Let $\delta=\frac{d-3}2$. By Corollary~\ref{cor:ball}, $c_{f_\beta}$ is a positive constant times
$J(k)/J(k+\delta)=\Gamma(k+1)\Gamma(k+\delta+\frac32)/(\Gamma(k+\frac32)\Gamma(k+\delta+1))$ with $k=k_\beta=\frac{d^2-1}2+d\beta$; as $k_\beta$ is strictly increasing in $\beta$, it suffices to find the sign of the
derivative in $k$ for $k\ge\frac{d^2-1}2$. Its logarithmic derivative in $k$ is $\psi(k+1)+\psi(k+\delta+\frac32)-\psi(k+\frac32)-\psi(k+\delta+1)$, with
$\psi=\Gamma'/\Gamma$ strictly concave on $(0,\infty)$ (as $\psi'(x)=\sum_{m\ge0}(x+m)^{-2}$ is strictly decreasing). The two pairs
of arguments have the same sum. If $d=3$ they coincide. If $d\ge4$, then $\delta\ge\frac12$ and
$k+1<k+\frac32\le k+\delta+1<k+\delta+\frac32$, so by strict concavity the derivative is negative. If $d=2$, then $\delta=-\frac12$, the
pairs are $\{k+1,k+1\}$ and $\{k+\frac32,k+\frac12\}$, and the derivative is positive.

(c) For a unit vector $\nu$ and $e\perp\nu$, the chord at offset $p\in[-1,1]$ is $\{p\nu+se:|s|\le w\}$, $w=(1-p^2)^{1/2}$, with $f=(1+2p^2+2s^2)/(2\pi)$ on it, and by
direct integration
\[
  g(p)=\int_{-w}^w\!\!\int_{-w}^w f(p\nu+se)f(p\nu+te)|s-t|\,ds\,dt=\frac{2}{35\pi^2}(1-p^2)^{3/2}(16p^4+52p^2+37),
\]
so $g'(p)=-\frac{2}{5\pi^2}p(1-p^2)^{1/2}(16p^4+28p^2+1)\le0$ on $[0,1]$. As $b_2=\frac12$ (proof of Corollary~\ref{cor:plane}),
Lemma~\ref{lem:ig}(c) gives $h_f(z)=\int_0^{2\pi}g(|z|\cos\theta)\,d\theta$, which is nonincreasing in $|z|$ since $g$ is even and
nonincreasing on $[0,1]$. So $m_f=\int_0^{2\pi}g(\cos\theta)\,d\theta$; with $\int_0^{2\pi}|\sin\theta|^3\cos^{2m}\theta\,d\theta=8/((2m+1)(2m+3))$ this
is $\frac{16}{35\pi^2}\bigl(\frac{16}{35}+\frac{52}{15}+\frac{37}3\bigr)=\frac{9104}{1225\pi^2}$, and $c_f=8/m_f=1225\pi^2/1138$. (The integrals
were also checked in exact arithmetic, \texttt{radial\_disk.py} in \cite{Repo}.)
\end{proof}

% ---- Proposition 7.3, strictly convex part, and Corollary 7.4 vertex value (lines 635-657, 673-674) ----
\begin{proposition}\label{prop:sym}
Let $d=2$ and let $\Omega$ be centrally symmetric; by isometry invariance we may take $\Omega=-\Omega$. For $z\in\partial\Omega$,
$h_\Omega(sz)$ is nonincreasing in $s\in[0,1]$, so $m_\Omega=\min_{\partial\Omega}h_\Omega$. If $\Omega$ is strictly convex, $h_\Omega(sz)$ is
strictly decreasing in $s\in[0,1]$, and all minimizers lie on $\partial\Omega$.
\end{proposition}

\begin{proof}[Proof of Proposition~\ref{prop:sym}]
Fix $\theta$, and let $\tau_\theta=\max_\Omega\ip\cdot{\nu_\theta}$. In the coordinates $q\nu_\theta+ye_\theta$, the chord at offset $q\in[-\tau_\theta,\tau_\theta]$ is
$\{y_-(q)\le y\le y_+(q)\}$ with $y_+$ concave and $y_-$ convex, so $L(q)=L(\theta,q)=y_+-y_-$ is concave, and even by symmetry;
hence nonincreasing in $|q|$. For $z\in\partial\Omega$, the chord through $sz$ in direction $e_\theta$ has offset $s\ip z{\nu_\theta}$, and
$|s\ip z{\nu_\theta}|\le \tau_\theta$; so by Corollary~\ref{cor:plane} $h_\Omega(sz)$ is nonincreasing in $s$. As $0\in\operatorname{int}\Omega$ (since $\Omega=-\Omega$), every point of $\Omega$
is $sz$ with $z\in\partial\Omega$, $s\in[0,1]$, so $m_\Omega=\min_{\partial\Omega}h_\Omega$. If $L(q)=L(q')$ with $0\le q<q'\le \tau_\theta$, concavity and monotonicity force
$L(0)=L(q')$, so $L$ is constant on $[-q',q']$, $y_+=y_-+L(0)$ is affine there, and $\partial\Omega$ contains the segment
$\{r\nu_\theta+y_+(r)e_\theta:|r|\le q'\}$. So if $\Omega$ is strictly convex, $L$ is strictly decreasing in $|q|$ on $[0,\tau_\theta]$. For
$0\le s<s'\le1$ and every $\theta$ with $\ip z{\nu_\theta}\ne0$ (all but one), $L(\theta,s\ip z{\nu_\theta})>L(\theta,s'\ip z{\nu_\theta})$, so
$h_\Omega(sz)>h_\Omega(s'z)$, and $h_\Omega>\min_{\partial\Omega}h_\Omega$ on $\operatorname{int}\Omega$.
\end{proof}

\begin{corollary}\label{cor:square}
For $\square=[0,1]^2$ and $\phi=(1+\sqrt5)/2$, the minimum of $h_\square$ over $\square$ is attained exactly at the midpoints of
the edges, $m_\square=\frac{\sqrt5}3+\frac{11}{12}\log\phi\approx1.1865$, and $c_\square=96/(4\sqrt5+11\log\phi)\approx6.743$. At the
vertices $h_\square=\frac23(\sqrt2+\log(1+\sqrt2))\approx1.5304$.
\end{corollary}
$A(\frac12)=\frac18\cdot\frac12(2\sqrt5+3\log\phi)+\frac12(\frac{\sqrt5}4+\log\phi)=\frac{\sqrt5}4+\frac{11}{16}\log\phi$. At a vertex,
$h_\square=\frac23A(1)$ and $A(1)=2\int_0^{\pi/4}\sec^3\theta\,d\theta=\sqrt2+\log(1+\sqrt2)$.

% ---- Remarks: Where is h smallest, Densities, Fluctuations (lines 686-704) ----
\section{Remarks}\label{sec:remarks}

\paragraph{Where is $h_\Omega$ smallest?} For hyperplanes, is $m_\Omega=\min_{\partial\Omega}h_\Omega$ for every convex body? For $\Omega=-\Omega$
in $d\ge3$, let $K_p=\{w\in\nu^\perp:w+p\nu\in\Omega\}$; then $K_{-p}=-K_p$, and by convexity $K_{(2s-1)p}\supseteq sK_p+(1-s)(-K_p)$ for
$s\in[0,1]$. So if $\Psi(sK+(1-s)(-K))\ge\Psi(K)$ for convex $K\subset\R^{d-1}$ and $s\in[\frac12,1]$ (for $d=2$, $K$ is a segment and
equality holds), then $\Psi(K_p)$ is nonincreasing in $|p|$, as $\Psi$ is monotone, and Lemma~\ref{lem:ig}(c) gives
$m_\Omega=\min_{\partial\Omega}h_\Omega$; we do not know whether this holds. For an equilateral and a long thin triangle, floating-point scans of
Corollary~\ref{cor:plane} also put the minimizers on the boundary (\texttt{h\_triangle\_scan.py}, \texttt{h\_triangle\_refine.py}
in \cite{Repo}).

\paragraph{Densities.} The density enters Theorem~\ref{thm:main} only through $h_f$, and it may vanish on part of $\Omega$, as the
spanned flats cross the whole body. Which continuous density on $\Omega$ minimizes $c_f$? For hyperplanes and the ball,
Proposition~\ref{prop:radial} shows that the uniform density is optimal among rotation-invariant densities in dimension $3$ only;
we do not know whether an optimal density exists for $d=2$ or $d\ge4$, nor whether it is rotation invariant.

\paragraph{Fluctuations.} For hyperplanes, by analogy with \cite{Janson86}, does $m_\Omega n^d\rho_\Omega^{(n)}/d!-d^2\log n-\beta'\log\log n$
converge in distribution to a shifted Gumbel law, and does the dependence through shared points affect the limit? We expect
$\beta'$ to depend on the set where $h_\Omega$ is minimal and on the growth of $h_\Omega$ near it, and so to differ between the disk
(minimal on the whole circle) and the square (four points).
